\section{Tổng quan ứng dụng và thiết kế cơ sở dữ liệu}
\label{sec:app-db-design}

% ===================================================================
\subsection{Tổng quan ứng dụng}
\label{subsec:app-overview}

\subsubsection{Giới thiệu ứng dụng}

\textbf{Novelle} là một ứng dụng web cho phép người dùng đọc và sáng tác truyện dài kỳ trực tuyến. Nội dung truyện do chính người dùng viết và đăng tải theo từng chương. Ứng dụng được xây dựng theo mô hình client--server: phía client là một ứng dụng web chạy trên trình duyệt, phía server cung cấp REST API để xử lý nghiệp vụ và truy xuất dữ liệu.

\subsubsection{Mục tiêu và phạm vi ứng dụng}

Ứng dụng hướng tới việc cung cấp đầy đủ các luồng chính của một nền tảng truyện trực tuyến, bao gồm: tìm và đọc truyện, quản lý thư viện cá nhân, sáng tác và phát hành chương, và quản trị nội dung. Phạm vi của ứng dụng giới hạn ở các chức năng cốt lõi này, trong đó ưu tiên các thao tác thêm, sửa, xoá và truy vấn dữ liệu được hiện thực đầy đủ.

\subsubsection{Các chức năng chính của ứng dụng}

\begin{itemize}
    \item \textbf{Trang chủ và tìm kiếm:} hiển thị truyện mới cập nhật, xem nhiều, đánh giá cao; tìm truyện theo từ khoá và lọc theo thể loại, thẻ, trạng thái, điểm đánh giá.
    \item \textbf{Chi tiết truyện và đọc chương:} xem thông tin truyện, mục lục, trang tác giả; đọc chương và tự động lưu vị trí đọc.
    \item \textbf{Thư viện cá nhân:} lịch sử đọc, truyện và tác giả đang theo dõi, tủ truyện tự tạo.
    \item \textbf{Đánh giá:} chấm điểm và viết nhận xét cho truyện.
    \item \textbf{Đăng ký làm tác giả:} gửi yêu cầu kèm bút danh và theo dõi trạng thái xét duyệt.
    \item \textbf{Studio sáng tác:} tạo và quản lý truyện, viết chương với bản nháp tự động lưu, chia tập, hẹn giờ đăng, xem lại các phiên bản chỉnh sửa, và xem thống kê lượt xem, lượt theo dõi, điểm đánh giá.
    \item \textbf{Quản trị:} quản lý và khoá tài khoản, duyệt yêu cầu tác giả, ẩn hoặc gỡ truyện và nhận xét vi phạm, quản lý thể loại và thẻ, xem thống kê toàn hệ thống.
\end{itemize}

% ===================================================================
\subsection{Yêu cầu chức năng của hệ thống}
\label{subsec:functional-requirements}

\subsubsection{Danh sách các yêu cầu chức năng}
% TODO

\subsubsection{Mô tả chi tiết các chức năng}
% TODO

% ===================================================================
\subsection{Phân tích và mô hình hóa dữ liệu}
\label{subsec:data-modeling}

\subsubsection{Xác định các thực thể dữ liệu}
% TODO

\subsubsection{Xác định thuộc tính của các thực thể}
% TODO

\subsubsection{Xác định mối quan hệ giữa các thực thể}
% TODO

\subsubsection{Xây dựng sơ đồ thực thể - liên kết (ERD)}
\label{subsubsec:erd}
% TODO: chèn sơ đồ ERD, ví dụ:
% \begin{figure}[H]
%     \centering
%     \includegraphics[width=\textwidth]{image/chapter3/erd.png}
%     \caption{Sơ đồ thực thể - liên kết của hệ thống}
%     \label{fig:erd}
% \end{figure}

% ===================================================================
\subsection{Thiết kế cơ sở dữ liệu trên MS SQL Server (SQL)}
\label{subsec:mssql-design}

\subsubsection{Xây dựng lược đồ quan hệ (Relational Schema)}
% TODO

\subsubsection{Thiết kế các bảng và thuộc tính}
% TODO

\subsubsection{Xác định khóa chính và khóa ngoại}
% TODO

\subsubsection{Xác định các ràng buộc dữ liệu}
% TODO

% ===================================================================
\subsection{Thiết kế dữ liệu trên Redis (NoSQL)}
\label{subsec:redis-design}

\subsubsection{Chiến lược tổ chức dữ liệu}
% TODO

\subsubsection{Thiết kế Key Naming Convention}
% TODO

\subsubsection{Lựa chọn cấu trúc dữ liệu Redis}
% TODO

\subsubsection{Thiết kế cách lưu trữ và truy xuất dữ liệu}
% TODO
